\section{Introduction}

\subsection{Background}
Embodied cognition proposes that mental processes emerge from the dynamic interaction between an organism's sensorimotor systems and its environment \citep{Glenberg2010, Barsalou2008}. Central to this framework is peripersonal space (PPS), defined as the region of space immediately surrounding the body within reaching distance \citep{rizzolatti1997space}, that acts as a multisensory interface that integrates tactile, proprioceptive, auditory, and visual signals to mediate interactions with nearby objects \citep{Serino2019, diPellegrino2015, serino2015bodypart}. Extensive research demonstrates that the boundaries of peripersonal space (PPS) are remarkably plastic, dynamically adjusting to a variety of factors and adapting flexibly to both action goals and contextual demands of cognition. These modulatory influences include tool use, which extends PPS boundaries to incorporate the tool \citep{iriki1996coding, maravita2004tools}; social context, which reconfigures PPS based on interpersonal dynamics \citep{teneggi2013social}; and experimental manipulations that alter self-location via body-swap illusions facilitated through virtual reality \citep{noel2015peripersonal}, as well as vestibular stimulation that modifies perceived spatial orientation of the body \citep{pfeiffer2018vestibular}. The dynamic flexibility of PPS representation underscores its function as an adaptive interface between the body and its immediate environment, supporting the organism's capacity for embedded, embodied forms of cognition. Emerging research further suggests that bodily states shape cognitive processes and perceptual frameworks by contributing to predictive models of the self and the world \citep{barsalou2008grounded, seth2013interoceptive, linson2018active, barrett2016theory}. Among these bodily signals, breathing stands out for its unique role in integrating interoceptive and exteroceptive cues—hence the aim of this study is to investigate how respiratory phase modulates peripersonal space. We briefly outline the role of breathing in mediating interoceptive and exteroceptive interactions.

\subsection{Breathing at the Interoceptive–Exteroceptive Interface}

Breathing, as a sensory modality, serves as a critical interface between internal bodily states and external environmental conditions. Unlike other autonomic functions, it is both voluntarily controllable and automatically driven by skeletal muscles, generating rich afferent signaling to the brain \citep{Betka2022}. \citet{Chan2024} describe respiratory interoception as a multi-layered process engaging various levels of the central nervous system for sensing, interpreting, integrating, and responding to information about internal bodily state. Afferent input arises from mechanoreceptors in the airways and lungs (responsive to stretch, flow, and irritation), chemoreceptors (monitoring blood gases), and proprioceptors in the respiratory muscles \citep{Varga2017}. These signals converge in brainstem nuclei and are relayed to subcortical and cortical areas, where two complementary systems shape the conscious experience of breathing: somatosensory regions encode physical attributes (e.g., location, timing, intensity), while limbic structures such as the amygdala and insula mediate affective evaluation \citep{Chan2024}. Respiratory interoception undergoes neural gating, where during normal breathing, sensory information remains below the threshold of conscious awareness. Only significant alterations in ventilatory patterns or deliberate attention can activate this regulatory gateway, permitting respiratory sensations to reach consciousness awareness through thalamic, hippocampal, and amygdala pathways. By contrast, the exteroceptive dimension of breathing remains relatively understudied, even though each breath brings us into contact with sensory receptors that detect temperature \citep{burgess1988effects, garefis2023human}, humidity \citep{Chan2024}, and olfactory information \citep{Mainland2006}—all of which convey information about the external environment. However, recently in cognitive psychology there has been a renewed interest in the role of breathing and other interoceptive signals in processing of external stimuli, such as emotional processing, as emotional stimuli are processed differently depending on their timing relative to the respiratory cycle \citep{Zelano2022}.

\section{Research Questions and Hypotheses}

This research investigates how respiratory phase modulates PPS boundaries. We propose that inhalation (characterized by increased lung volume) expands PPS representation, while exhalation (decreased lung volume) contracts PPS representation. This hypothesis aligns with embodied cognition theories suggesting that bodily expansion and contraction correspondingly affect cognitive representations of numerosity \citep{belli2021breathink}.

\subsection{Rationale and Design}
Study 1 aims to establish the basic effects of respiratory phase on PPS processing using a controlled box breathing paradigm. Specifically, we will employ a within-subjects design to test for differences in peripersonal space boundaries between respiratory phases during breath holds following inhalation versus exhalation. In Study 2, we repeat this procedure; however, instead of wearing blindfolds, participants will wear a VR headset that provides exteroceptive feedback about their interoceptive breathing state. This method—termed \textit{exteroceptive–interoceptive sensory substitution} \citep{Goral2024}—has been applied in various immersive virtual reality and art installations that use breath as a bioresponsive feedback modality to create environments responsive to users' respiration. Most previous studies investigating this phenomenon have relied solely on subjective measures. Here, we aim to assess its effects on the embodiment of space using a more robust design that incorporates objective measures.
